% ============ cover.tex --- 封面与"关于本重排版"页 ============
% 由 AJbook.cls 的 coverpage 选项(coverpage=cover)引入
% 设计仿自《代数学方法》封面风格(李文威)，图饰采用微局部分析相空间与波前集特征流
% 原书: Peter Hintz, An Introduction to Microlocal Analysis (Springer GTM 304, 2025)

\setCJKfamilyfont{coverfont}{SimHei}	% 设置书名字体(Fandol 环境下回退)
\setCJKfamilyfont{cover-author-font}{SimHei}
\colorlet{octa}{cyan!50!gray}
\definecolor{ANTnavy}{HTML}{16223B}
\definecolor{ANTgold}{HTML}{B08D57}
\definecolor{ANTgoldlight}{HTML}{D9C29B}
\definecolor{ANTgrey}{HTML}{8C93A6}

\begin{titlepage}
\begin{tikzpicture}[remember picture, overlay]

  % ---------- outer frame ----------
  \draw[ANTnavy, line width=1.1pt]
       ($(current page.north west)+(1.4cm,-1.4cm)$) rectangle
       ($(current page.south east)+(-1.4cm,1.4cm)$);
  \draw[ANTgold, line width=0.4pt]
       ($(current page.north west)+(1.6cm,-1.6cm)$) rectangle
       ($(current page.south east)+(-1.6cm,1.6cm)$);

  % ---------- corner ornaments ----------
  \begin{scope}[shift={($(current page.north west)+(1.4cm,-1.4cm)$)}]
    \draw[ANTgold, line width=0.9pt] (0,0) -- (0.55,0);
    \draw[ANTgold, line width=0.9pt] (0,0) -- (0,-0.55);
    \filldraw[ANTgold] (0,0) circle (1.6pt);
  \end{scope}
  \begin{scope}[shift={($(current page.north east)+(-1.4cm,-1.4cm)$)}]
    \draw[ANTgold, line width=0.9pt] (0,0) -- (-0.55,0);
    \draw[ANTgold, line width=0.9pt] (0,0) -- (0,-0.55);
    \filldraw[ANTgold] (0,0) circle (1.6pt);
  \end{scope}
  \begin{scope}[shift={($(current page.south west)+(1.4cm,1.4cm)$)}]
    \draw[ANTgold, line width=0.9pt] (0,0) -- (0.55,0);
    \draw[ANTgold, line width=0.9pt] (0,0) -- (0,0.55);
    \filldraw[ANTgold] (0,0) circle (1.6pt);
  \end{scope}
  \begin{scope}[shift={($(current page.south east)+(-1.4cm,1.4cm)$)}]
    \draw[ANTgold, line width=0.9pt] (0,0) -- (-0.55,0);
    \draw[ANTgold, line width=0.9pt] (0,0) -- (0,0.55);
    \filldraw[ANTgold] (0,0) circle (1.6pt);
  \end{scope}

  % ---------- top ornament: 相空间余切球丛与奇性传播流线示意 ----------
  \begin{scope}[shift={($(current page.north)+(0,-5.5cm)$)}]
    % 底流形 M
    \draw[ANTgrey, line width=0.8pt] (-3.2,-1.2) .. controls (-1,-0.8) and (1,-1.5) .. (3.2,-1.1);
    \node[ANTnavy, font=\scriptsize] at (3.5,-1.1) {$M$};
    
    % 底流形上的奇性支集 supp u 与奇异点 x_0
    \filldraw[ANTgold] (0,-1.05) circle (1.8pt) node[below=2pt, font=\scriptsize, text=ANTnavy] {$x_0$};
    
    % 余切纤维 T^*_{x_0} M \setminus \{0\}
    \draw[ANTnavy, dashed, line width=0.6pt] (0,-0.8) -- (0,1.8);
    \node[ANTnavy, font=\scriptsize] at (-1.2,1.8) {$T^*_{x_0} M \setminus \{0\}$};
    
    % 特征锥面 / 特征集 Char(P)
    \draw[ANTgrey!60, line width=0.7pt] (-2.2,1.3) -- (0,-1.05) -- (2.2,1.3);
    \draw[ANTgrey!50, line width=0.5pt] (0,1.3) ellipse (2.2cm and 0.45cm);
    
    % 波前集中的方向 (x_0, \xi_0) \in WF(u)
    \draw[ANTgold, line width=1.2pt, -stealth] (0,-1.05) -- (1.2,0.6) node[midway, right=2pt, font=\scriptsize] {$\xi_0$};
    \filldraw[ANTgold] (1.2,0.6) circle (2pt) node[right=3pt, font=\scriptsize, text=ANTnavy] {$(x_0,\xi_0) \in \mathrm{WF}(u)$};
    
    % Hamilton 向量场流线 / 双曲零测地线
    \draw[ANTnavy, line width=1.1pt, -{Stealth[length=2.5mm]}]
         (-2.0,0.2) .. controls (-0.5,0.7) and (0.8,0.5) .. (2.3,1.3)
         node[above right=1pt, font=\scriptsize] {$\exp(t H_p)(x_0,\xi_0)$};
    
    \node[ANTnavy, font=\scriptsize\itshape, align=center] at (0,-1.9)
         {微局部分析: 相空间中的几何与分析\,\textbullet\,波前集与奇性沿特征流线的传播};
  \end{scope}

  % ---------- title block ----------
  \node[ANTgold, font=\scriptsize\sffamily] at ($(current page.center)+(0,2.35cm)$)
       {GRADUATE TEXTS IN MATHEMATICS \quad 304};

  \draw[ANTgold, line width=0.5pt] ($(current page.center)+(-3.6cm,1.95cm)$) -- ($(current page.center)+(3.6cm,1.95cm)$);

  \node[ANTnavy, align=center] at ($(current page.center)+(0,0.75cm)$)
       {\Huge\bfseries 微局部分析导论\\[0.45em] \large An Introduction to Microlocal Analysis};

  \draw[ANTgold, line width=0.5pt] ($(current page.center)+(-3.6cm,-0.75cm)$) -- ($(current page.center)+(3.6cm,-0.75cm)$);

  \node[ANTnavy, font=\normalsize\itshape] at ($(current page.center)+(0,-1.55cm)$)
       {伪微分算子\,\textbullet\,波前集\,\textbullet\,奇性传播\,\textbullet\,双曲与波动方程};

  % ---------- bottom : author / date / formula ----------
  \node[ANTnavy, font=\large] at ($(current page.south)+(0,5.0cm)$) {Peter Hintz\ 著};
  \node[ANTgrey, font=\normalsize] at ($(current page.south)+(0,4.2cm)$) {Springer Nature Switzerland AG, 2025};

  \node[ANTgold, font=\small] at ($(current page.south)+(0,3.15cm)$)
       {$\displaystyle \mathrm{WF}(u) \setminus \mathrm{WF}(Pu) \subset \operatorname{Char}(P) \subset T^*M \setminus 0$};

  % ---------- translator credit ----------
  \node[ANTnavy, font=\large] at ($(current page.south)+(0,2.15cm)$)
       {翻译及重排: ZCode};

\end{tikzpicture}
\vfill

\end{titlepage}

% ---------- 封面背面: 关于本重排版 ----------
\thispagestyle{empty}
\begin{center}
    \small
    \noindent\emph{关于本重排版}\\[0.6em]
    \parbox{0.86\textwidth}{
    本文件为 Peter Hintz 教授（苏黎世联邦理工学院 ETH Zürich）所著
    \emph{An Introduction to Microlocal Analysis}
    （Springer Graduate Texts in Mathematics, Volume 304, 2025 年出版，ISBN: 978-3-031-90705-0）
    的中文非官方 \LaTeX{} 翻译与重排版。
    \medskip

    重排版严格遵循原书结构与数学推导脉络：
    \begin{itemize}[leftmargin=2em, itemsep=0.2em]
      \item 排版系统采用李文威《代数学方法》开源模板体系（\texttt{AJbook.cls}）定制，使用 \hologo{XeLaTeX} + \texttt{Biber} + \texttt{texindy} 编译；
      \item 全书定理类环境与原书 GTM 体例严格保持一致，按章统排；
      \item 全书 83 条参考文献与原书书末 Bibliography 保持一对一精确映射与双向跳转；
      \item 术语翻译参考国内微局部分析与现代偏微分方程学术界通用定名；
      \item 译本仅供学术研究与教学交流使用，原书版权属于原作者及 Springer Nature Switzerland AG。
    \end{itemize}
    }
\end{center}
\vfill
\clearpage
